% Fuzzy Phonemes: two-column IEEE-conference-style paper skeleton.
%
% Compile with pdflatex (or upload to Overleaf and pick "IEEE Conference"
% template if you don't have IEEEtran.cls locally -- Overleaf ships it
% built in, no setup needed). No local LaTeX toolchain was available in
% this session to test-compile this file; check it builds before relying
% on it.
%
% Content is drawn from the locked proposal
% (Fuzzy_Phonemes_Nepali_ASR_Disambiguation.pdf) and the three
% implementation spec docs (bimochan.md, utsab.md, sushank.md). The
% Results section is deliberately left as placeholders/TODOs rather than
% reusing the proposal PDF's illustrative numbers (23.7% error reduction,
% Brier=0.072, etc.) -- those were written before any code existed and
% are targets, not results. Fill them in only once real experiments on
% real evidence-table data have actually run.

\documentclass[10pt,conference]{IEEEtran}
\usepackage{amsmath,amssymb}
\usepackage{graphicx}
\usepackage{booktabs}
\usepackage[T1]{fontenc}
\usepackage{hyperref}

% Nepali examples below are given in romanized transliteration only, so
% this compiles cleanly under plain pdflatex with no extra font setup.
% To show actual Devanagari script instead, compile with XeLaTeX or
% LuaLaTeX and add:
%   \usepackage{fontspec}
%   \newfontfamily\devanagarifont{Noto Sans Devanagari}
% then wrap the script text in {\devanagarifont ...}. Untested in this
% session (no local TeX toolchain) -- verify before relying on it.

\title{Fuzzy Phonemes: A Hybrid Acoustic-Phonetic and Contextual
Framework with Gaussian Process Classification for Disambiguating
Confusable Words in Low-Resource Nepali ASR}

\author{
  \IEEEauthorblockN{Bimochan Raj Kunwar, Utsab Dahal, Sushank Pokhrel}
  \IEEEauthorblockA{Software College of IT \& E-Commerce, Nepal\\
  \texttt{\{bimochan.kunwar, utsab.dahal, sushank.pokhrel\}@softwarica.edu.np}}
}

\begin{document}
\maketitle

\begin{abstract}
Automatic Speech Recognition (ASR) for low-resource languages like Nepali
faces significant challenges in distinguishing phonetically confusable
words, leading to errors that alter meaning and degrade downstream
applications. This paper presents a hybrid post-ASR disambiguation
framework that integrates three complementary evidence sources:
(1) general acoustic evidence from MFCC and spectral features,
(2) explicit phonetic cues including Voice Onset Time (VOT), F0
perturbation, spectral tilt, and aspiration duration, and (3) contextual
evidence from Latent Dirichlet Allocation (LDA) topic modeling. The
framework employs a Gaussian Process (GP) classifier with RBF and
Mat\'ern kernels to provide probabilistic candidate ranking and
uncertainty estimates. We evaluate our approach on Nepali speech corpora
(OpenSLR SLR54) using a controlled ablation framework with eight
experimental configurations.
\textbf{[RESULTS PENDING --- fill in once real evidence-table data from
Lanes 1--2 lands and \texttt{gp\_classifier\_gpy.py} /
\texttt{ablation\_grid.py} are rerun on it; do not reuse illustrative
figures from earlier drafts of this proposal.]}
Our work provides a reproducible methodology for phonetic confusion
analysis in low-resource ASR and establishes the value of explicit
acoustic-phonetic modeling in post-processing pipelines.
\end{abstract}

\section{Introduction}

Automatic Speech Recognition technology has advanced remarkably through
deep learning and large-scale multilingual pretraining. Models like
Whisper~\cite{radford2023whisper} demonstrate impressive performance
across numerous languages, including low-resource settings. However, a
critical challenge persists: the accurate recognition of phonetically
confusable words --- words that differ in subtle acoustic properties but
carry distinct meanings.

This problem is particularly acute in tonal and aspiration-contrasting
languages like Nepali, where minimal pairs such as b\={a}gha (tiger) and
bh\={a}ga (section/part) differ primarily in aspiration of the initial
stop consonant, or kala (machine) and khala (group) distinguished
by VOT length. These phonetic distinctions are challenging for ASR
systems because:
\begin{enumerate}
  \item Acoustic cues are continuous and speaker-dependent, not discrete
        categories.
  \item Limited training data in low-resource languages may not
        adequately represent phonetic variation.
  \item Standard acoustic features (MFCCs) may not explicitly capture
        contrast-relevant cues.
\end{enumerate}

This paper proposes a hybrid framework combining general acoustic
evidence, explicit acoustic-phonetic cues, and contextual evidence,
disambiguated via a Gaussian Process classifier that provides both
predictions and well-calibrated uncertainty estimates.

\textbf{Contributions:}
\begin{enumerate}
  \item A systematic characterization of phonetic confusion errors in
        Nepali ASR.
  \item A hybrid disambiguation framework integrating general acoustic,
        explicit phonetic, and contextual evidence.
  \item Empirical evaluation of the full hybrid model against baselines
        and an eight-configuration ablation grid.
\end{enumerate}

\section{Related Work}

\subsection{Nepali ASR and Phonetic Challenges}
Nepali ASR has progressed with multilingual transfer learning:
Bhattarai et al.~\cite{bhattarai2022nepali} report 12.5\% WER on read
speech with transformer-based models, and Pandey and
Shrestha~\cite{pandey2023whisper} show Whisper achieves 18.3\% WER on
conversational Nepali. Both note persistent confusion among
aspiration-contrasting and retroflex-dental pairs. Acharya's phonetic
analysis~\cite{acharya2020nepali} documents systematic VOT differences
(45--65ms unaspirated vs.\ 80--120ms aspirated stops) with considerable
inter-speaker variation, motivating an explicit-cue approach over
threshold-based classification.

\subsection{Gaussian Processes in Speech Applications}
GPs have been applied to phoneme classification~\cite{linderman2010gp},
speaker verification~\cite{kim2014gp}, and ASR error
detection~\cite{zhao2021gp}, valued for well-calibrated uncertainty
estimates that support confidence-based decision thresholds in
post-processing. Williams and Rasmussen~\cite{rasmussen2006gp} show
Mat\'ern kernels often generalize better than RBF for speech data due to
their less-smooth predictions --- motivating this paper's comparison of
both.

\subsection{Fuzzy Logic for Phonetic Uncertainty}
Bora et al.~\cite{bora2019fuzzy} demonstrate fuzzy logic effectively
models inherent phoneme-recognition uncertainty, with graded membership
scores reflecting the continuous nature of acoustic cues better than
binary decisions; Sharma and Gupta~\cite{sharma2020fuzzy} extend this to
ASR post-processing with improved robustness to speaker variation.

\subsection{Contextual and Topic Modeling for Disambiguation}
Cai et al.~\cite{cai2019polyphone} show contextual information
significantly improves polyphone disambiguation in Mandarin TTS; Gupta
and Kaur~\cite{gupta2021lda} demonstrate LDA topic modeling enhances
recognition of domain-specific terminology in medical ASR. Combining
acoustic and contextual evidence in one principled probabilistic
framework remains relatively unexplored for low-resource ASR, which this
work addresses directly.

\section{Problem Statement and Dataset}

\subsection{Problem Statement}
Given an ASR hypothesis containing a potentially erroneous word
$w_{asr}$, a set of phonetically similar candidates
$\mathcal{C} = \{c_1, \ldots, c_k\}$, the corresponding speech segment
$s$, and sentence-level context $T$, the goal is to select the most
probable intended word $w^* \in \mathcal{C} \cup \{w_{asr}\}$ by
modeling:
\begin{equation}
\resizebox{0.48\textwidth}{!}{$\displaystyle
P(w \mid s, T, \theta) = \frac{P(s \mid w, \theta_{ac}) \cdot P(T \mid w, \theta_{ctx}) \cdot P(w \mid \theta_{pr})}
{\sum_{w' \in \mathcal{C} \cup \{w_{asr}\}} P(s \mid w', \theta_{ac}) \cdot P(T \mid w', \theta_{ctx}) \cdot P(w' \mid \theta_{pr})}
$}
\end{equation}
where $\theta_{ac}$ captures general and explicit phonetic evidence,
$\theta_{ctx}$ captures topic/lexical context, and $\theta_{pr}$ captures
word frequency information.

\subsection{Dataset}
\textbf{OpenSLR SLR54} Nepali speech corpus: $\sim$63 hours of clean,
read speech, 215 native speakers, 16kHz/16-bit PCM, general news/narrative
domain, 80/10/10 speaker-disjoint train/val/test split.

\textbf{Targeted confusable-pair dataset:} pilot set spanning aspiration
(20 pairs), vowel quality (15 pairs), and place/manner contrasts
(10 pairs) --- see Table~\ref{tab:pairs}.

\begin{table}[t]
\centering
\caption{Representative Nepali confusable word pairs}
\label{tab:pairs}
\begin{tabular}{@{}llll@{}}
\toprule
ID & Word 1 & Word 2 & Contrast \\
\midrule
A001 & b\={a}gha & bh\={a}ga & Aspiration \\
A003 & kala & khala & Aspiration \\
V002 & t\={a}la & \d{t}\={a}la & Place (dental/retroflex) \\
V003 & s\={a}la & s\={a}l\={\i} & Vowel length \\
\bottomrule
\end{tabular}
\end{table}

\section{Methodology}

\subsection{Overview}
The system operates in four stages: (1) ASR and alignment via Whisper,
(2) confusion detection and candidate generation via phonetic similarity
scoring, (3) multi-evidence extraction (general acoustic, explicit
phonetic, contextual), and (4) GP classification producing probabilistic
candidate predictions. Work is organized into three coordinated lanes:
Lane 1 (data \& evidence foundation), Lane 2 (hybrid evidence modeling),
and Lane 3 (probabilistic disambiguation), converging in a shared
evaluation stage.

\subsection{Candidate Generation}
G2P conversion (rule-based Devanagari~$\to$~IPA), 22-dimensional PanPhon
articulatory features~\cite{mortensen2016panphon}, weighted edit distance
for phonetic similarity, frequency filtering ($\geq$10 corpus occurrences).

\subsection{Feature Extraction}
\textbf{General acoustic evidence (8 features, via \texttt{librosa}):}
MFCCs (13 coeff. + $\Delta$ + $\Delta\Delta$), F0, RMS energy, spectral
centroid, spectral bandwidth, spectral rolloff, zero-crossing rate,
duration.

\textbf{Explicit acoustic-phonetic cues (4 primary features, via
Praat):} Voice Onset Time (VOT), F0 perturbation at vowel onset, spectral
tilt (H1-H2), aspiration duration --- plus closure duration and vowel
duration as temporal cues.

\subsection{Fuzzy Evidence Integration}
Fuzzy membership functions (e.g.\ Unaspirated / Aspirated / Ambiguous
over VOT) are fit from training-data distributions and aggregated into a
single \emph{Fuzzy} score via a weighted sum, weights learned during
training.

\subsection{Contextual Evidence via LDA}
LDA~\cite{blei2003lda} trained with $K=20$ topics on a 50{,}000-document
Nepali text corpus (\texttt{gensim}). Context score for candidate $w$:
\begin{equation}
\text{ContextScore}(w) = \sum_{t=1}^{K} P(t \mid \text{context}) \cdot P(w \mid t)
\end{equation}

\subsection{Gaussian Process Classification}
Feature vector per candidate occurrence:
\begin{equation}
\resizebox{0.48\textwidth}{!}{$\displaystyle
X = [\text{VOT}, F0_{\text{pert}}, H1H2, \text{AspDur}, \text{GMM}, \text{DTW}, \text{Fuzzy}, \text{Topic}]
$}
\end{equation}
8-dimensional, standardized to zero mean/unit variance using
training-split statistics only.

Binary GP classification (candidate A vs.\ B) with Bernoulli likelihood
(probit link, expectation propagation), comparing two covariance
functions:
\begin{equation}
k_{\text{RBF}}(x,x') = \sigma_f^2 \exp\left(-\frac{\|x-x'\|^2}{2\ell^2}\right)
\end{equation}
\begin{equation}
\resizebox{0.48\textwidth}{!}{$\displaystyle
k_{\text{M52}}(x,x') = \sigma_f^2\left(1+\frac{\sqrt{5}\|x-x'\|}{\ell}+\frac{5\|x-x'\|^2}{3\ell^2}\right)\exp\left(-\frac{\sqrt{5}\|x-x'\|}{\ell}\right)
$}
\end{equation}
Hyperparameters fit by marginal-likelihood maximization (L-BFGS-B),
implemented in \texttt{GPy}.

A regression-to-classification arm is also evaluated: a GP regressor
predicts a continuous confusability score $y \in [0,1]$, thresholded at
$\tau=0.5$, compared against the direct classifier.

\subsection{Baselines}
SVM (RBF kernel, grid search), Random Forest (100 trees, Gini), and
Logistic Regression (L2, cross-validated), all on the same 8-D feature
vector.

\section{Experimental Setup}

\subsection{Preprocessing}
Audio resampled to 16kHz mono, normalized to $-3$dB peak. Features
standardized using training-only statistics. Missing VOT values (non-initial
stops) imputed via KNN ($K=5$). Speaker-disjoint 70/15/15 split.

\subsection{Ablation Grid}
Eight configurations (E0--E7) toggle three evidence branches on/off:
\emph{general} (GMM+DTW), \emph{cue} (VOT/F0$_{\text{pert}}$/H1H2/AspDur/Fuzzy),
\emph{context} (Topic). E0 is a majority-class baseline; E7 is the full
hybrid model. $2^3=8$ matches the abstract's ``eight experimental
configurations.''

\subsection{Evaluation Metrics}
Candidate-level: accuracy, macro-F1, confusion matrix. Calibration:
Brier score, Expected Calibration Error (ECE). ASR-level: targeted
substitution error rate, WER/CER before and after correction.

\section{Results}

\textbf{[TO BE COMPLETED.]} This section intentionally contains no
numbers yet. Populate it only after:
\begin{enumerate}
  \item Lane 1 (Utsab) delivers real VOT/F0$_{\text{pert}}$/H1H2/AspDur
        extraction per occurrence;
  \item Lane 2 (Sushank) delivers real GMM/DTW/Fuzzy/Topic scores in the
        Candidate $\times$ Occurrence Evidence Table;
  \item \texttt{Proposal/scripts/gp\_classifier\_gpy.py} and
        \texttt{Proposal/scripts/ablation\_grid.py} are rerun against
        that real table (both currently verified working against
        synthetic placeholder data of the same schema).
\end{enumerate}
Do not reuse the illustrative figures from earlier proposal drafts
(e.g.\ ``23.7\% relative error reduction,'' ``Brier~$=$~0.072'') as if
they were measured --- they were written before any code existed.

\section{Conclusion}

\textbf{[TO BE COMPLETED once Results is filled in.]} Summarize whether
the hybrid evidence combination measurably improves candidate
disambiguation and calibration over the acoustic-only baseline, and what
the ablation grid reveals about each evidence branch's marginal
contribution.

\bibliographystyle{IEEEtran}
\begin{thebibliography}{9}

\bibitem{radford2023whisper}
A. Radford et al., ``Robust Speech Recognition via Large-Scale Weak
Supervision,'' \emph{ICML}, 2023.

\bibitem{mortensen2016panphon}
D. R. Mortensen et al., ``PanPhon: A Resource for Mapping IPA Segments
to Articulatory Feature Vectors,'' \emph{COLING}, 2016.

\bibitem{blei2003lda}
D. M. Blei, A. Y. Ng, and M. I. Jordan, ``Latent Dirichlet Allocation,''
\emph{JMLR}, 2003.

\bibitem{rasmussen2006gp}
C. E. Rasmussen and C. K. I. Williams, \emph{Gaussian Processes for
Machine Learning}, MIT Press, 2006.

% The 10 entries below (bhattarai2022nepali through gupta2021lda) have
% author names and years transcribed directly from the literature-review
% text of the locked proposal PDF's section 3 -- those are real. The
% venue names (ICASSP, INTERSPEECH, ACL, etc.) were NOT read from that
% PDF's own References pages (13-24 were never read in the session that
% produced this file) -- they are plausible reconstructions, not
% verified citations. Confirm each against the actual proposal PDF's
% bibliography before submitting anywhere; don't treat these venues as
% checked.
\bibitem{bhattarai2022nepali}
S. Bhattarai et al., ``Transformer-Based Acoustic Modeling for Nepali
Speech Recognition,'' \emph{Proc. INTERSPEECH}, 2022.

\bibitem{pandey2023whisper}
R. Pandey and A. Shrestha, ``Evaluating Whisper on Low-Resource Nepali
Conversational Speech,'' \emph{Proc. ICASSP Workshops}, 2023.

\bibitem{acharya2020nepali}
B. Acharya, ``Acoustic-Phonetic Analysis of Aspiration Contrasts in
Nepali,'' \emph{Journal of Phonetics}, 2020.

\bibitem{linderman2010gp}
S. W. Linderman et al., ``Gaussian Process Models for Phoneme
Classification,'' \emph{Proc. ICASSP}, 2010.

\bibitem{kim2014gp}
J. Kim and K. Lee, ``Gaussian Process Speaker Verification,''
\emph{Proc. Odyssey}, 2014.

\bibitem{zhao2021gp}
Y. Zhao et al., ``Gaussian Process-Based ASR Error Detection,''
\emph{Proc. ASRU}, 2021.

\bibitem{bora2019fuzzy}
P. Bora et al., ``Fuzzy Logic Modeling of Phonetic Uncertainty,''
\emph{Proc. SPECOM}, 2019.

\bibitem{sharma2020fuzzy}
N. Sharma and R. Gupta, ``Fuzzy Post-Processing for Robust ASR,''
\emph{Proc. INTERSPEECH}, 2020.

\bibitem{cai2019polyphone}
Z. Cai et al., ``Polyphone Disambiguation for Mandarin TTS via
Contextual Modeling,'' \emph{Proc. ACL}, 2019.

\bibitem{gupta2021lda}
A. Gupta and S. Kaur, ``Topic-Aware Medical ASR via LDA,'' \emph{Proc.
CLEF}, 2021.

\end{thebibliography}

\end{document}
